% !PRINTED 56
\weberpage{56}

11. \gesperrt{Zähler und Nenner dieser Darstellung sind durch $\varphi(u)$ selbst, abgesehen von Einheitsfaktoren, eindeutig bestimmt.}

Die Funktionen $G(u)$, $G_1(u)$ können in Primfaktoren, d.\,h. in Faktoren zerlegt werden, die nur in einem Punkte Null werden. Verfolgt man diesen Weg bei den doppelt periodischen Funktionen, so gelangt man zu den \gesperrt{Weierstrassschen} $\sigma$-Funktionen. Wir werden aber einen anderen Weg gehen, auf dem wir den $\sigma$-Funktionen erst an einer späteren Stelle begegnen.

\begin{center}
\textbf{\S\,16. Periodizit\"at.}
\end{center}

Eine Funktion $\varphi(u)$ von $u$, welche die Eigenschaft hat, daß für ein konstantes $\omega$ und für jeden Wert von $u$
\weq{1}{\varphi(u+\omega) = \varphi(u)}
ist, heißt \gesperrt{periodisch} und $\omega$ heißt die \gesperrt{Periode}.

Der Gegenstand unserer Untersuchung sind Funktionen $\varphi(u)$ mit zwei Perioden $\omega_1, \omega_2$, von denen wir ein für allemal voraussetzen wollen, daß ihr Verhältnis $\omega_2:\omega_1$ nicht reell und daß der imaginäre Teil dieses Verhältnisses \gesperrt{positiv} sei. In der letzteren Annahme liegt nur eine Festsetzung über die Bezeichnung $\omega_1, \omega_2$. Denn wenn $\omega_2:\omega_1$ einen negativen imaginären Teil hat, so ist der von $\omega_1:\omega_2$ positiv.

Jede Kombination $m_1\omega_1 + m_2\omega_2$ ist, wenn $m_1, m_2$ ganze Zahlen sind, gleichfalls eine Periode.

Wir nehmen in der Ebene $u$ einen beliebigen Punkt $u_0$ an und bezeichnen die vier Punkte
\[
u_0, \quad u_0 + \omega_1, \quad u_0 + \omega_1 + \omega_2, \quad u_0 + \omega_2;
\]
verbinden wir diese vier Punkte der Reihe nach durch gerade Linien, indem wir vom letzten zum ersten zurückkehren, so erhalten wir ein Parallelogramm, welches das \gesperrt{Periodenparallelogramm} genannt wird. Es können die geradlinigen Seiten des Periodenparallelogramms auch durch krummlinige Züge ersetzt werden, wenn nur die gegenüberliegenden durch Parallelverschiebung zur Deckung gelangen und die einzelnen Züge keine Schleifen bilden.

Ist $\omega_1 = \alpha_1 + \beta_1 i$, $\omega_2 = \alpha_2 + \beta_3 i$ und $\alpha_1, \alpha_2, \beta_1, \beta_3$ reell, % [sic] Printed beta_3 twice; mathematically expected beta_2.
so ist nach der Voraussetzung $(\alpha_1\beta_2 - \alpha_2\beta_1)$ positiv, und die geometrische Bedeutung dieser Größe ist (nach bekannten Sätzen
%CONT>>
